\textbf{(CAP-COLUNI 2020 -- Q14)} Em janeiro de 2017, um mesmo modelo de tênis era vendido em duas lojas (A e B) pelo mesmo preço. Ao final do primeiro ano de vendas, a loja A aumentou o preço dessa mercadoria em 10\%, enquanto a loja B manteve o preço inalterado. Ao final do segundo ano, o preço dessa mercadoria foi novamente aumentado pela loja A em 10\%, já na loja B foi dado um desconto de 5\%. Ao final dos dois anos de vendas, os preços desse tênis nas lojas A e B ($P_A$ e $P_B$), respectivamente, eram diferentes.

Com base nessas informações, o percentual de reajuste que deve ser aplicado a $P_B$ para se obter igualdade de valores com $P_A$ é de, aproximadamente:

\begin{enumerate}
	\item 20,6\%
	\item 22,3\%
	\item 25,9\%
	\item 27,4\%
\end{enumerate}
